% =============================================================================
\chapter{Meshes, Reference Cells, and Geometric Mappings}
\label{ch:meshes}
% =============================================================================

This chapter introduces the geometric framework behind finite element
discretization. It covers mesh topology, reference-cell constructions, and the
maps that transfer basis functions, quadrature rules, and differential
operators between physical and reference domains.

\section{Cells, facets, edges, and mesh topology}
\label{sec:meshes:topology}

% Simplicial and tensor-product meshes, entity incidence, and orientation

\section{Admissibility, regularity, and mesh quality}
\label{sec:meshes:quality}

% Shape regularity, quasi-uniformity, and geometric assumptions in analysis

\section{Reference cells and local coordinates}
\label{sec:meshes:reference_cells}

% Canonical intervals, triangles, tetrahedra, quadrilaterals, and hexahedra

\section{Affine, isoparametric, and curvilinear mappings}
\label{sec:meshes:mappings}

% Pullbacks from physical cells to reference cells and approximation of geometry

\section{Jacobians, metrics, and transformed derivatives}
\label{sec:meshes:jacobians}

% Determinants, inverse Jacobians, metric tensors, and geometric scaling

\section{Mesh refinement and adaptivity concepts}
\label{sec:meshes:refinement}

% h-refinement, anisotropy, hanging nodes, and geometric multilevel structure
